%%
%% Whitepaper Shell / Briefing Outline for Black Hat USA 2027 (Briefings / Arsenal)
%%
\documentclass[11pt,a4paper]{article}
\usepackage[utf8]{inputenc}
\usepackage{booktabs}
\usepackage{hyperref}
\usepackage{microtype}
\usepackage{geometry}
\geometry{margin=1in}

\title{\textbf{Trust Me, I'm an Issue: Weaponising and Remediating Prompt-Injection Exposure in CI Coding Agents}}

\author{
  \textbf{Author / Presenter Team}\\
  \texttt{research@example.invalid}
}
\date{Black Hat USA 2027 --- Briefings \& Arsenal Tracks}

\begin{document}

\maketitle

\begin{abstract}
AI coding agents deployed across continuous integration (CI) workflows are transforming repository automation, but they introduce a severe Confused Deputy vulnerability: delegating write-capable forge tokens to unauthenticated Web actors. When a repository wires autonomous agents to open webhook events---such as public issues, pull requests, and discussions---without strict perimeter gating, any outsider can inject prompt payloads that compel the agent to execute arbitrary tools. In this briefing and Arsenal tool demonstration, we unveil CI-Trigger-Scan, an automated open-source security scanner designed to detect vulnerable agent workflow wiring. Drawing on an empirical census of 367 agent repositories and 448 workflows, we show that 67.0\% execute on untrusted Web triggers and 88.3\% grant write permissions, with 95.8\% omitting maintainer identity checks. We demonstrate live exploit mechanics: weaponising public issues to force coding agents to dump CI secrets, modify production source branches, and compromise downstream releases. We dissect how default starter templates published by AI vendors (such as Anthropic Claude Code) propagate ungated configurations to 86.8\% of downstream adoptions, contrasting them with secure-by-default patterns. Finally, we demonstrate practical defenses: implementing identity-verified perimeter gating and split-token sandboxing that neutralises 96.0\% of attack paths while preserving developer utility.
\end{abstract}

\section{Introduction: Autonomous Coding Agents in the CI Pipeline}
Continuous integration (CI) runners were originally architected for deterministic build and test pipelines. With the arrival of autonomous LLM coding agents, maintainers are hooking agentic workflows directly to forge events (issues, comments, pull requests) to automatically triage bugs, draft documentation, and generate code patches. This architecture inverts traditional trust assumptions: unauthenticated Web text directly drives code execution in privileged runners.

\section{Vulnerability Analysis: Confused Deputy on Code Forges}
A public forge issue is inherently untrusted. However, when an issue triggers an agent workflow, the forge injects an authenticated token (\texttt{GITHUB\_TOKEN}) with elevated permissions into the runner. If the agent parses the issue content as system-level instructions rather than passive data, an indirect prompt injection attack enables the external adversary to commandeer the privileged token.

\section{Empirical Exposure Census}
We present findings from our census of 367 agent-enabled repositories across 448 GitHub Actions workflows:
\begin{itemize}
    \item \textbf{Untrusted Triggers:} 67.0\% (300/448) execute on untrusted events (\texttt{issues}, \texttt{pull\_request}, \texttt{pull\_request\_target}).
    \item \textbf{Write Authority:} 88.3\% (416/471) of agent jobs operate with write permissions (\texttt{contents:write}, \texttt{pull-requests:write}).
    \item \textbf{Perimeter Gating Omission:} 95.8\% of untrusted-trigger workflows lack identity filters (\texttt{github.actor} or \texttt{author\_association}).
    \item \textbf{Exploitable Exposure:} 64.2\% of surveyed repositories grant write tokens directly to unauthenticated issue creators.
\end{itemize}

\section{Live Exploitation Demonstration}
We demonstrate step-by-step weaponisation of ungated CI agent workflows:
\begin{enumerate}
    \item \textbf{Payload Delivery:} Filing an issue with a payload that mimics natural bug reports while embedding adversarial instructions.
    \item \textbf{Token Hijacking:} Instructing the agent to inspect environment variables and exfiltrate secrets via outbound HTTP requests or issue comments.
    \item \textbf{Supply Chain Compromise:} Directing the agent to commit backdoored commits into repository branches or tag malicious releases.
\end{enumerate}

\section{The Upstream Cause: Vendor Starter Templates}
Our investigation reveals that insecure defaults in vendor documentation directly cause downstream exposure:
\begin{itemize}
    \item Documentation templates omitting perimeter gates (e.g., Anthropic Claude Code) resulted in 86.8\% of real-world adoptions remaining fully exposed.
    \item Templates enforcing identity gating by default (e.g., Google Gemini CLI) achieved 100\% secure adoptions in the wild.
\end{itemize}

\section{Tool Demonstration: CI-Trigger-Scan}
At Black Hat Arsenal, we present \texttt{CI-Trigger-Scan}, a lightweight, zero-dependency static analysis tool that parses GitHub Actions workflow ASTs, extracts trigger-permission tuples, evaluates perimeter conditions, and produces actionable risk scorecards for repository maintainers.

\section{Defense-in-Depth Remediations}
We detail robust, production-ready remediations:
\begin{itemize}
    \item \textbf{Perimeter Gating:} Enforcing strict \texttt{author\_association} checks prior to runner dispatch.
    \item \textbf{Split-Token Sandboxing:} Restricting reasoning and patch generation to read-only environments, delegating write actions to separate human-in-the-loop approval gates.
    \item \textbf{Tool Capability Containment:} Restricting agent tool calls to ephemeral containers without ambient credential exposure.
\end{itemize}

\section{Conclusion}
Treating CI agent workflows as trusted environments creates critical supply-chain vulnerabilities. CI-Trigger-Scan empowers the security community and repository maintainers to identify and eliminate these confused deputy exposures before adversaries exploit them.

\end{document}
